% !TeX program = xelatex
\documentclass[10.5pt,a4paper]{article}

\usepackage[margin=1.35cm]{geometry}
\usepackage{titlesec}
\usepackage{enumitem}
\usepackage[hidelinks]{hyperref}
\usepackage{tabularx}
\usepackage{array}
\usepackage{xcolor}

\setlength{\parindent}{0pt}
\setlength{\parskip}{0pt}
\pagestyle{empty}

\titleformat{\section}{\normalsize\bfseries}{}{0pt}{}[\vspace{-0.55em}\rule{\textwidth}{0.4pt}]
\titlespacing*{\section}{0pt}{0.55em}{0.28em}
\setlist[itemize]{leftmargin=1.15em, itemsep=0.08em, topsep=0.08em, parsep=0pt}

\newcommand{\resumeSubheading}[4]{%
  \begin{tabularx}{\textwidth}{@{}X r@{}}
    \textbf{#1} & #2 \\
    \textit{#3} & \textit{#4} \\
  \end{tabularx}\vspace{-0.15em}
}

\newcommand{\resumeExperience}[4]{%
  \begin{tabularx}{\textwidth}{@{}X r@{}}
    \textbf{#1} & #2 \\
    {\small #3} & {\small #4} \\
  \end{tabularx}\vspace{0.10em}
}

\newcommand{\resumeProject}[3]{%
  \begin{tabularx}{\textwidth}{@{}X r@{}}
    \textbf{#1} & #2 \\
  \end{tabularx}
  {\small\textit{#3}}\vspace{0.12em}
}

\newcommand{\resumeParagraph}[1]{%
  {\setlength{\parskip}{0pt}\linespread{1.08}\selectfont #1\par}\vspace{0.10em}
}

\begin{document}

\begin{center}
    {\LARGE \textbf{Fubin Dou}} \\
    \vspace{0.25em}
    Email: \href{mailto:doufb@mail.ustc.edu.cn}{doufb@mail.ustc.edu.cn}
    \quad | \quad
    Phone: +86 15522881185
\end{center}
\vspace{-0.6em}

\section{Education}
\resumeSubheading
{University of Science and Technology of China, School of Mathematical Sciences}{Anhui, China}
{B.S. in Computational Mathematics}{Sept. 2023 -- Jun. 2027 (Expected)}

\section{Honors and Awards}
\begin{itemize}
    \item Outstanding Student Scholarship, 2025
\end{itemize}

\section{Research Interests}
\resumeParagraph{My research interests lie at the intersection of applied mathematics, artificial intelligence, computer vision, and computer graphics, with a focus on dynamic 3D reconstruction, neural rendering, physics-constrained deformable modeling, optimization and inverse problems, scientific computing, and AI for Science.}

\section{Publications}
\begin{itemize}
    \item Haoqin Hong*, Ding Fan*, \textbf{Fubin Dou}, Zhi-Li Zhou, Haoran Sun, Congcong Zhu, Jingrun Chen.
    \textit{Physics-Informed Deformable Gaussian Splatting: Towards Unified Constitutive Laws for Time-Evolving Material Field}.
    Association for the Advancement of Artificial Intelligence (AAAI), 2026. \textbf{Accepted}.
\end{itemize}

\section{Research Experience}
\resumeExperience
{SCAI Lab, University of Science and Technology of China}{Anhui, China \quad Nov. 2024 -- Present}
{Advisor: Prof. Jingrun Chen, Department of Mathematics, USTC}{}
\resumeParagraph{Participated in research on a physics-informed deformable 3D Gaussian Splatting framework that introduces continuum mechanics into dynamic scene reconstruction. I worked on the static--dynamic decoupling and dynamic mask modules, separating moving foreground regions from static backgrounds to alleviate camera/object motion ambiguity and improve both static background geometry and dynamic foreground motion consistency.}

\section{Academic Projects}
\resumeProject
{Differentiable Physics-Driven Spring-Mass System and Inverse Reconstruction of Dynamic Scenes}{University of Science and Technology of China}
{Project on deformable dynamic scene modeling and inverse reconstruction}
\resumeParagraph{Built a differentiable physics simulation framework based on spring-mass systems and embedded the simulation process into an optimization loop to infer scene states and material parameters from observed dynamic sequences. The project connects discrete spring systems with continuum mechanics, fits Neo-Hookean elastic energy to a discrete topology through energy matching, estimates physical parameters such as Young's modulus and Poisson's ratio, and implements simulation and automatic differentiation workflows with Taichi. It further uses differentiable rendering and physics losses to optimize dynamic deformation, providing an optimization-based approach for recovering geometry, motion, and physical properties of deformable objects from visual observations.}

\section{Relevant Coursework}
\begin{tabularx}{\textwidth}{>{\raggedright\arraybackslash}X >{\raggedright\arraybackslash}X >{\raggedright\arraybackslash}X}
Real Analysis (4.0) & Complex Analysis (4.0) & Functional Analysis (4.0) \\
Multivariable Calculus (4.0) & Mathematical Statistics (4.0) & Numerical Algebra (4.0)
\end{tabularx}

\section{Skills}
\begin{itemize}
    \item \textbf{Programming Languages:} C, C++, Python, \LaTeX
    \item \textbf{Deep Learning and Scientific Computing:} PyTorch, Taichi, NumPy
    \item \textbf{Tools and Platforms:} Linux, Git
\end{itemize}

\end{document}
